\begin{lem}[One-step $\flowsto$ Soundness]
	\label{lem:flowsto_soundness_emptytrace}
	For every $ c \in \cdom, D \in \dbdom$ and two assignment events $\event_1, \event_2 \in \eventset^{\asn}$,
	if $\eventdep(\event_1, \event_2, [\event_1; \event_2],  c, D) $,
	then there is a flows to relation, $\flowsto(\pi_1(\event_1)^{\pi_2(\event_1)}, \pi_1(\event_2)^{\pi_2(\event_2)}, c)$.
	%
	\[
	\begin{array}{l}
		\forall \event_1, \event_2 \in \eventset^{\asn}, c \in \cdom, D \in \dbdom 
		\sthat  
		\eventdep(\event_1, \event_2, [\event_1; \event_2],  c, D) 
		\\ \quad 
		\implies 
		\flowsto(\pi_1(\event_1)^{\pi_2(\event_1)}, \pi_1(\event_2)^{\pi_2(\event_2)}, c)
	\end{array}
	\]
\end{lem}
Proof Summary:
\\
1. case of (the labelled unique assignment command associated to the $\event_2$ 
is executed but the value assigned to the variable in this event is changed in second execution)
\\
show x directly used by the assignment of the second event
\\
2.(the labelled unique assignment command associated to the $\event_2$ isn't executed in second execution)
\\
show x is directly used by the boolean expression for a conditional command and second event shows in the body of that conditional command 
%
\begin{proof}
	By the Definition~\ref{def:event_dep} for $\eventdep(\event_1, \event_2, [\event_1; \event_2], c, D)$, 
	we know there are 2 cases:
	%
	\caseL{Value Dependency in Equation~\ref{eq:eventdep}(b)}
		% \textbf{(the labelled unique assignment command associated to the $\event_2$ 
		% is executed but the value assigned to the variable in this event is changed in second execution).}
\input{appendices/thm-adaptivity/base_case_val}
	% \end{case}
	\caseL{Control Dependency in Equation~\ref{eq:eventdep}(c)}
		% \textbf{(the labelled unique assignment command associated to the $\event_2$ isn't executed in second execution).}
		\\
		Proof Summary:
		\\
		1. Let $\event_b$ be the testing event,
		in the same way of case 1, we get:
		 $\pi_1(\event_1) \in \kw{FV}(\pi_1(\event_b)) 
		 \land 
		%  l_1 \in \live^{l_b}$
		\pi_1(\event_1)^{l_1} \in \live(l_b, c)$
		%
		 \\
		 2. By Lemma~\ref{lem:ctldep_inv}, we know:
		   $\forall z \in \kw{FV}(\pi_1(\event_b)) \sthat  \exists i \in \mathbb{N} \sthat 
		 \flowsto(z^i, \pi_1(\event)^{\pi_2(\event)}, c)$
		 %
		 \\
		 3. By $\flowsto$ definition we have:
		   $\flowsto(\pi_1(\event_1)^{\pi_2(\event_1)}, \pi_1(\event_2)^{\pi_2(\event_2)}, c)$
		\input{appendices/thm-adaptivity/base_case_ctl}
\end{proof}
%


\begin{lem}[Multiple-step $\flowsto$ Soundness]
	\label{lem:depevents_exist}
For every $D \in \dbdom , c \in \cdom, \trace \in \ftdom$, and two assignment events 
$\event_1, \event_2 \in \eventset^{\asn}$,
if the trace $trace$ has the form $\trace = [\event_1] \tracecat \trace' \tracecat [\event_2]$ with $\trace' \in \ftdom$, and 
if $\eventdep(\event_1, \event_2, \trace, c, D)$,
then either there is a flows to relation between the labeled variables from $\event_1$ and $\event_2$, i.e.,
$\flowsto(\pi_1(\event_1)^{\pi_2(\event_1)}, \pi_1(\event_2)^{\pi_2(\event_2)}, c) $,
or 
there exists
$\event \in \trace'$ such that
$
\eventdep(\event_1, \event, \trace[\event_1:\event], c, D)
$
and $\flowsto(\pi_1(\event)^{\pi_2(\event)}, \pi_1(\event_2)^{\pi_2(\event_2)}, c)$ if $\event \in \eventset^{\asn}$
or 
$\flowsto(z^l, \pi_1(\event_2)^{\pi_2(\event_2)}, c)$ if $\event \in \eventset^{\test}$
for every $z \in \kw{FV}(\pi_2(\event))$ and some label $l \in \lvar(c)$.
%
	\[
	\begin{array}{l}
		\forall D \in \dbdom , c \in \cdom, \trace \in \ftdom \sthat  \forall \event_1, \event_2 \in \eventset^{\asn} \sthat 
		 \exists \trace' \in \ftdom \sthat  
		 \\ \quad 
		 \trace = [\event_1] \tracecat \trace' \tracecat [\event_2]
		\land
		\eventdep(\event_1, \event_2, \trace, c, D) 
		% \\ \quad 
		\implies 
		\flowsto(\pi_1(\event_1)^{\pi_2(\event_1)}, \pi_1(\event_2)^{\pi_2(\event_2)}, c) 
		\\ \qquad  \lor
		\Big( 		
		\exists \event \in \trace' \sthat  
		\eventdep(\event_1, \event, \trace[\event_1:\event], c, D)
		\\ \qquad \quad
		\land 
		\event \in \eventset^{\asn} \implies \flowsto(\pi_1(\event)^{\pi_2(\event)}, \pi_1(\event_2)^{\pi_2(\event_2)}, c)
		\\ \qquad \quad
		\land 
		\event \in \eventset^{\test} \implies \forall z \in \kw{FV}(\pi_2(\event)) \sthat \exists l \in \lvar(c) \sthat
		\flowsto(z^l, \pi_1(\event_2)^{\pi_2(\event_2)}, c)
		\Big) 
		% \\ \qquad \quad  \lor
		% \Big( 		
		% \exists \event \in \eventset^{\test} \sthat  
		% \event \in \trace' \land
		% 	\eventdep(\event_1, \event, \trace[\event_1:\event], c, D)
		% \land 
		% \flowsto(\pi_1(\event)^{\pi_2(\event)}, \pi_1(\event_2)^{\pi_2(\event_2)}, c) 
		% \Big) 
		% \\ \qquad \qquad \lor
		% \flowsto(\pi_1(\event_1)^{\pi_2(\event_1)}, \pi_1(\event_2)^{\pi_2(\event_2)}, c) 
	\end{array}
	\]
\end{lem}
\begin{proof}
	Taking arbitrary 
	$ D \in \dbdom , c \in \cdom, \trace \in \ftdom $ and two events 
	$\event_1, \event_2 \in \eventset^{\asn}$, where $\trace$ has the form 
	$\trace = [\event_1] \tracecat \trace' \tracecat [\event_2]$ 
	for some 
	$\trace' \in \ftdom$ and $\eventdep(\event_1, \event_2, \trace, c, D)$.
	\\
	Then it is sufficient to show
	\[
		\begin{array}{l}
			\neg\Big( 		
			\exists \event \in \trace' \sthat  
			\eventdep(\event_1, \event, \trace[\event_1:\event], c, D)
			\\ \qquad 
			\land 
			\big( 
				\event \in \eventset^{\asn} \implies 
				\flowsto(\pi_1(\event)^{\pi_2(\event)}, \pi_1(\event_2)^{\pi_2(\event_2)}, c)
			\big)
			\\ \qquad
			\land 
			\big(
			\event \in \eventset^{\test} \implies \forall z \in \kw{FV}(\pi_2(\event)) \sthat \exists l \in \lvar(c) \sthat
			\flowsto(z^l, \pi_1(\event_2)^{\pi_2(\event_2)}, c)
			\big)
			\Big) 
			\\ \qquad \quad
			\implies
			\flowsto(\pi_1(\event)^{\pi_2(\event)}, \pi_1(\event_2)^{\pi_2(\event_2)}, c)
		\end{array}
		\]
		Take arbitrary $ \event \in \trace'$, it is sufficient to show
		\[
			\begin{array}{l}
				\Big( 		
				\neg \eventdep(\event_1, \event, \trace[\event_1:\event], c, D) ~(\star)
				\\ \qquad 
				\lor 
				\neg \big( 
					\event \in \eventset^{\asn} \implies 
					\flowsto(\pi_1(\event)^{\pi_2(\event)}, \pi_1(\event_2)^{\pi_2(\event_2)}, c)
				\big) ~(\diamond)
				\\ \qquad
				\lor 
				\neg \big(
				\event \in \eventset^{\test} \implies \forall z \in \kw{FV}(\pi_2(\event)) \sthat \exists l \in \lvar(c) \sthat
				\flowsto(z^l, \pi_1(\event_2)^{\pi_2(\event_2)}, c)
				\big)
				\Big)  ~(\square)
				\\ \qquad \quad
				\implies
				\flowsto(\pi_1(\event)^{\pi_2(\event)}, \pi_1(\event_2)^{\pi_2(\event_2)}, c)
			\end{array}
		\]
		Splitting the preconditions, it is sufficient to show the following cases.
		\[
			\begin{array}{l}
			(1) (\star) \implies \flowsto(\pi_1(\event)^{\pi_2(\event)}, \pi_1(\event_2)^{\pi_2(\event_2)}, c)
			\\
			(2) 
			\neg(\star) \land (\diamond) \implies \flowsto(\pi_1(\event)^{\pi_2(\event)}, \pi_1(\event_2)^{\pi_2(\event_2)}, c)
			\\
			(3)
			\neg(\star) \land (\square) \implies \flowsto(\pi_1(\event)^{\pi_2(\event)}, \pi_1(\event_2)^{\pi_2(\event_2)}, c)
		\end{array}
		\]
	\caseL{$(\star) \implies \flowsto(\pi_1(\event)^{\pi_2(\event)}, \pi_1(\event_2)^{\pi_2(\event_2)}, c)$}
	This case is proved by the first case in Lemma~\ref{lem:inv_indepevents_asn}.

	\caseL{$\neg(\star) \land (\diamond) \implies \flowsto(\pi_1(\event)^{\pi_2(\event)}, \pi_1(\event_2)^{\pi_2(\event_2)}, c)$}
	This case is proved by the second case in Lemma~\ref{lem:inv_indepevents_asn}.

	\caseL{$\neg(\star) \land (\square) \implies \flowsto(\pi_1(\event)^{\pi_2(\event)}, \pi_1(\event_2)^{\pi_2(\event_2)}, c)$}
	This case is proved by Lemma~\ref{lem:inv_indepevents_ctl} and Lemma~\ref{lem:ctldep_inv}. 
\end{proof}
%
%
%
%
\begin{lem}[Multi-step $\flowsto$ through Assignment Event]
		\label{lem:inv_indepevents_asn}
		For every $D \in \dbdom , c \in \cdom, \trace \in \ftdom$, one assignment events 
		$\event_1\in \eventset^{\asn}$, and another event $\event_2 \in \eventset$,
		if the trace $\trace$ has the form $\trace = [\event_1] \tracecat \trace' \tracecat [\event_2]$ with $\trace' \in \ftdom$, 
		and $\eventdep(\event_1, \event_2, \trace, c, D)$,
		then the following two conclusions hold when $\event_2$ is an assignment.
		%  event and a testing event respectively.
	% \begin{itemize}
	% 	\item
		If $\event_2 \in \eventset^{\asn}$,
		% For every $D \in \dbdom , c \in \cdom, \trace \in \ftdom$, and two assignment events 
		% $\event_1, \event_2 \in \eventset^{\asn}$,
		% if the trace $\trace$ has the form $\trace = [\event_1] \tracecat \trace' \tracecat [\event_2]$ with $\trace' \in \ftdom$, and $\eventdep(\event_1, \event_2, \trace, c, D)$,
		then for every $\event \in \trace'$, if it either doesn't have the \emph{Event May-Dependency} relation on $\event_1$, 
		or $\pi_1(\event)^{\pi_2(\event)}$ doesn't have the $\flowsto$ relation with $ \pi_1(\event_2)^{\pi_2(\event_2)}$,
		then the labelled variable $\pi_1(\event_1)^{\pi_2(\event_1)}$ directly flows to the other one $\pi_1(\event_2)^{\pi_2(\event_2)}$, 
		i.e., $\flowsto(\pi_1(\event_1)^{\pi_2(\event_1)}, \pi_1(\event_2)^{\pi_2(\event_2)}, c)$.
		%
		\[
		\begin{array}{l}
			\forall D \in \dbdom , c \in \cdom, \trace \in \ftdom \sthat  \forall \event_1, \event_2 \in \eventset^{\asn} \sthat 
			 \exists \trace' \in \ftdom \sthat
			 \\ \quad 
			 \trace = [\event_1] \tracecat \trace' \tracecat [\event_2]
			 \land
			\eventdep(\event_1, \event_2, \trace, c, D) 
			\\ \quad 
			\land 
			\Big( 
			\forall \event \in \trace' \sthat  
				\neg \eventdep(\event_1, \event, \trace[\event_1:\event], c, D) ~(\star)
				\\ \qquad 
				\lor 
				\neg \big( 
					\event \in \eventset^{\asn} \implies 
					\flowsto(\pi_1(\event)^{\pi_2(\event)}, \pi_1(\event_2)^{\pi_2(\event_2)}, c)
				\big) ~(\diamond)
			\Big) 
			\\ \quad 
			\implies 
			\flowsto(\pi_1(\event_1)^{\pi_2(\event_1)}, \pi_1(\event_2)^{\pi_2(\event_2)}, c)
		\end{array}
		\]
% 		\item
% If $\event_2 \in \eventset^{\test}$, 
% then for every $\event \in \trace'$, if it either doesn't have the \emph{Event May-Dependency} relations on $\event_1$,
% or $\pi_1(\event) \notin \kw{FV}(\pi_1(\event_2)) $,
% then 
% $\pi_1(\event_1) \in \kw{FV}(\pi_1(\event_2))$, and $ {\pi_2(\event_1)} = \llabel(\trace)$
% %
% \[
% \begin{array}{l}
% 	\forall D \in \dbdom , c \in \cdom, \trace \in \ftdom \sthat  \forall \event_1,\in \eventset^{\asn}, \event_2 \in \eventset^{\test} \sthat 
% 	 \exists \trace' \in \ftdom \sthat
% 	 \\ \quad 
% 	 \trace = [\event_1] \tracecat \trace' \tracecat [\event_2]
% 	 \land
% 	\eventdep(\event_1, \event_2, \trace, c, D) 
% 	\\ \quad 
% 	\land 
% 	\Big( \nexists \event \in \trace' \sthat  
% 	\pi_1(\event) \in \kw{FV}(\pi_1(\event_2)) \land
% 	\eventdep(\event_1, \event, \trace[\event_1:\event], c, D)
% 	\Big) 
% 	\\ \quad 
% 	\implies 
% 	\pi_1(\event_1) \in \kw{FV}(\pi_1(\event_2)) \land {\pi_2(\event_1)} = \llabel(\trace)
% \end{array}
% \]
% \end{itemize}
\end{lem}
%
\begin{proof}
Taking arbitrary $D \in \dbdom , c \in \cdom$, and an assignment events $\event_1 \in \eventset^{\asn}$ and another event 
$\event_2\in \eventset$.
\\
Without loss of generalization, 
taking arbitrary trace has the form $\trace = [\event_1, \ldots, \event_2]$ 
% arbitrary $\trace_2 \in \ftdom$,
%  then we know 
 and $\trace_2 \in \ftdom$ such that
%   $ $ and
 \[
	\trace = [\event_1] \tracecat \trace_2 \tracecat [\event_2]
	\land
	\eventdep(\event_1, \event_2, \trace, c, D)
	\]
% 	Then we have the following two cases by type of $\event_2$.
% %  let $\trace_2$ be this $\trace'$.
% %
% \caseL{$\event_2 \in \eventset^{\asn}$}
%
By the definition of $\eventdep(\event_1, \event_2, \trace, c, D)$ in Definition~\ref{def:event_dep}, 
taking $ \event_1', \event_2' \in \eventset^{\asn},
\trace_2' \in \ftdom, c_1, c_2 \in \cdom$ as the events, traces and commands
such that
% %  satisfying the definition,
%  we have following two executions:
% \[
%   \exists \event_1', \event_2' \in \eventset^{\asn},
%   \trace_2' \in \ftdom, c_1, c_2 \in \cdom \sthat 
% \]
%
\[
\begin{array}{l}
\config{c, \trace_0} \rightarrow^{*}
\config{c_1, \trace_1 \tracecat [\event_1]} \rightarrow^{*} \config{c_2, \trace_1 \tracecat [\event_1] \tracecat \trace_2 \tracecat [\event_2]} 
\\ \land
% \land
\config{c_1, \trace_1 \tracecat [\event_1']} \rightarrow^{*} \config{c_2, \trace_1 \tracecat [\event_1'] \tracecat \trace_2' \tracecat [\event_2']} 
\\
\land \diff(\event_2,\event_2' ) ~(\star) 
\\
\land 
    \vcounter(\trace_2, \pi_2(\event_2))
    = 
    \vcounter(\trace_2', \pi_2(\event_2'))
\end{array}
\]
%
%
By inversion Lemma.~\ref{lem:inv_event} on $\event_2$ and $\event_2'$ in the two executions
and $\diff(\event_2, \event_2')$,
 we have the following two execution instances:
\[
	\begin{array}{l}
		\config{c_1, \trace_1 \tracecat [\event_1]} \rightarrow^{*} \config{[\assign{\pi_1(\event_2)}{\expr_2 / \query(\qexpr_2)}]{}^{\pi_2(\event_2)};c_2, \trace_1 \tracecat [\event_1] \tracecat \trace_2} 
		\\ \quad
		\rightarrow^\rname{assn / query} \config{c_2, \trace_1 \tracecat [\event_1] \tracecat \trace_2 \tracecat [\event_2]}  
	\end{array}
	\]
%
\[
	\begin{array}{l}
		\config{c_1, \trace_1 \tracecat [\event_1']} \rightarrow^{*} \config{[\assign{\pi_1(\event_2)}{\expr_2 / \query(\qexpr_2)}]{}^{\pi_2(\event_2)};c_2, \trace_1 \tracecat [\event_1'] \tracecat \trace_2'} 
		\\ \quad
		\rightarrow^\rname{assn / query} \config{c_2, \trace_1 \tracecat [\event_1'] \tracecat \trace_2' \tracecat [\event_2']} 
	\end{array}
\]
, where $\expr_2 / \qexpr_2$ is the expression of the assignment command associated to the $\event_2$ and $\event_2'$ by the Inversion Lemma.~\ref{lem:inv_event}.
\\
Taking arbitrary $\event_z \in \trace_2$ and $\event_z \in \eventset^{\asn}$, by the assumption we know 
\[
\neg \eventdep(\event_1, \event, \trace[\event_1:\event_z], c, D)
\lor  \neg \flowsto(\pi_1(\event_z)^{\pi_2(\event_z)}, \pi_1(\event_2)^{\pi_2(\event_2)}, c).
\]
Splitting the formula above, we have the following two subcases.
\subcaseL{$\neg \eventdep(\event_1, \event, \trace[\event_1:\event_z], c, D)$}
In case of $\neg \eventdep(\event_1, \event, \trace[\event_1:\event_z], c, D)$,
by Definition~\ref{def:event_dep}, we know 
$\event_z$ shows up identically on $\trace_2'$ on the same position, i.e.,
% \[
% 	\nexists \event_z' \in \trace_2' \shat
% 	 \diff(\event_z, \event_z')
% \]
% $\event_z \in \trace_2'$ 
Then by Definition~\ref{def:adapt_env}, we know 
\[
	\env(\trace_1 \tracecat [\event_1] \tracecat \trace_2[:\event_z]) \pi_1(\event_z) = \env(\trace_1 \tracecat [\event_1'] \tracecat \trace_2'[:\event_z]) \pi_1(\event_z)
\]
Then we have 
\[
	\forall \event_z \in \trace_2 \sthat 
	\env(\trace_1 \tracecat \trace[\event_1:\event_z]) \pi_1(\event_z) = \env(\trace_1 \tracecat [\event_1'] \tracecat \trace_2'[:\event_z]) \pi_1(\event_z)
\]
Then by Definition~\ref{def:adapt_env}, we know 
\[
	\forall \event_z \in \trace_2 \sthat 
	\env(\trace_1 \tracecat [\event_1] \tracecat \trace_2) \pi_1(\event_z) = \env(\trace_1 \tracecat [\event_1] \tracecat \trace_2) \pi_1(\event_z)
\]
Then by the determinism of the execution, we know $\event_2 = \event_2'$.
Then we have the following by Definition~\ref{def:event_eq}.
\[
	\Big( (\pi_4(\event_2) = \bullet \land \pi_4(\event_2') = \bullet) \land \pi_3(\event_2) = \pi_3(\event_2') \Big)
	\lor
		\Big((\pi_4(\event_2) \neq \bullet \land \pi_4(\event_2') \neq \bullet) \land \pi_4(\event_2) = \pi_4(\event_2') \Big)
		\]
Then by Definition~\ref{def:diff}, we know
\[
	\neg \diff(\event_2, \event_2').
	\]
We also know $\diff(\event_2, \event_2')$ by $(\star)$, and we have this case proved by the contradiction.
%
\subcaseL{$ \eventdep(\event_1, \event, \trace[\event_1:\event_z], c, D) \land \neg \flowsto(\pi_1(\event_z)^{\pi_2(\event_z)}, \pi_1(\event_2)^{\pi_2(\event_2)}, c)$}
In this case, by Definition~\ref{def:feasible_flowsto}, we know
\[
	\pi_1(\event_z) \notin \kw{FV}(\expr_2 / \qexpr_2).
	\]
% In case of $ \pi_1(\event_z) \notin \kw{FV}(\expr_2 / \qexpr_2)$, by 
By Inversion Lemma~\ref{lem:inv_expr}
 of arithmetic and query expression cases, we know:
%
\[
	\begin{array}{l}
		\forall x^i \in \mathcal{LV}, \trace, \trace' \in \ftdom, v, v' \sthat 
		\\ \qquad 
		\Big( \forall z^j \in \mathcal{LV} \setminus \{\pi_1(\event_z)^{\pi_2(\event_z)} \} \sthat  
		\\ \qquad 
		\env(\trace) z = \env(\trace') z \Big) \land 
	\config{\trace, \expr_2 / \qexpr_2} \aarrow v \land \config{\trace', \expr_2} \aarrow  v' \\ \qquad 
	\implies v = v'
\end{array}
\]
	\[
		\begin{array}{l}
			\forall x^i \in \mathcal{LV}, \trace, \trace' \in \ftdom, \qval, \qval' \sthat 
			\\ \qquad 
		\Big( \forall z^j \in \mathcal{LV} \setminus \{\pi_1(\event_z)^{\pi_2(\event_z)} \} \sthat  
		\\ \qquad 
		\env(\trace) z = \env(\trace') z \Big) \land 
		\config{\trace, \qexpr_2} \qarrow \qval \land \config{\trace', \qexpr_2} \qarrow \qval' 
		\\ \qquad 
		\implies \qval =_{q} \qval'
	\end{array}
	\]
for $\expr_2$ or $\qexpr_2$ respectively.
\\
Let $\kw{use}_{\trace_2}$ a subset of the events in $\trace_2$, satisfying: 
\[
	\begin{array}{l}
		\forall \event \in \eventset^{\asn} \sthat  
	\event \in \kw{use}_{\trace_2} \Longleftrightarrow 
	\event \in \trace_2 \land
	\pi_1(\event) \in \kw{FV}(\expr_2 / \qexpr_2)
\end{array}		
\]
Let $\lvar({\trace_2}) $ denotes the set of all labeled variables in $\trace_2$
and $\lvar({\kw{use}_{\trace_2}})$ the set of all labeled variables in ${\kw{use}_{\trace_2}}$.

Then we also known for every $\event_z \in \kw{use}_{\trace_2}$, 
$\neg \eventdep(\event_1, \event_z, \trace[\event_1:\event_z], c, D)$
as follows.
\[
	\forall z^l \in \lvar(c) \setminus 
	\big( 
		( \lvar({\trace_2}) \setminus \lvar({\kw{use}_{\trace_2}})) \cup \{\pi_1(\event_1)^{\pi_2(\event_1)}\} \big)
	\sthat 
	\env(\trace_1 \tracecat [\event_1] \tracecat \trace_2) z = \env(\trace_1 \tracecat [\event_1'] \tracecat \trace_2') z
	~ (1)
\]
and 
\[
	\begin{array}{l}
	\forall z^l \in \lvar(c) \setminus ( \lvar({\trace_2}) \setminus \lvar({\kw{use}_{\trace_2}})), 
	\trace, \trace' \in \ftdom, v, v' \sthat  
	\\ \qquad
	\env(\trace) z = \env(\trace') z 
	\land 
	\config{\trace, \expr_2} \aarrow v 
	\land 
	\config{\trace', \expr_2} \aarrow v'
	\implies 
	v = v' 
	~ (2a)
\end{array}
\]
and
\[
	\begin{array}{l}
	\forall z^l \in \lvar(c) \setminus ( \lvar({\trace_2}) \setminus \lvar({\kw{use}_{\trace_2}})), 
	\trace, \trace' \in \ftdom, \qval, \qval' \sthat  
	\\ \qquad
	\env(\trace) z = \env(\trace') z 
	\land 
	\config{\trace, \qexpr_2} \qarrow \qval 
	\land 
	\config{\trace', \qexpr_2} \qarrow \qval'
	\implies 
	\qval =_q \qval' 
	~ (2q)
\end{array}
\]
%
Since $\diff(\event_2, \event_2')$, we also know:
%
\[	
\config{\trace_1 \tracecat [\event_1] \tracecat \trace_2, \expr_2} \aarrow \pi_3(\event_2)
\land 
\config{\trace_1 \tracecat [\event_1'] \tracecat \trace_2', \expr_2} \aarrow \pi_3(\event_2') 
\land 
\pi_3(\event_2) \neq \pi_3(\event_2')
\]
We know $\event_1$ is the only cause of the difference in $\event_y$ and $\event_y'$ when evaluating
the command 
$\clabel{\assign{\pi_1(\event_2)}{\expr_2 / \query(\qexpr_2)}}{}^{\pi_2(\event_2)}$.
%
\\
By inversion Lemma.~\ref{lem:inv_expr_gnl} of arithmetic and query expression cases, 
given the two traces
$\trace_1 \tracecat [\event_1'] \tracecat \trace_2'$ and 
$\trace_1 \tracecat [\event_1'] \tracecat \trace_2'$ 
satisfying this lemma by $(1)$, $(2a)$ and $(2q)$, 
we know
\[
  \pi_1(\event_1) \in \kw{FV}(\expr_2 / \qexpr_2) \land {\pi_2(\event_1)} = \llabel(\trace_1 \tracecat [\event_1] \tracecat \trace_2) \pi_1(\event_1) 
\]
%
By $\flowsto$ definition:
% \todo{add the Liveness inversion}
\[
  \flowsto(\pi_1(\event_1)^{\pi_2(\event_1)}, \pi_1(\event_y)^{\pi_2(\event_y)}, c)
\]
This case is proved.
%
\end{proof}
\begin{lem}[Multi-step $\flowsto$ Soundness through Testing Event]
	\label{lem:inv_indepevents_ctl}
	For every $D \in \dbdom , c \in \cdom, \trace \in \ftdom$, one assignment events 
	$\event_1\in \eventset^{\asn}$, and another event $\event_2 \in \eventset$,
	if the trace $\trace$ has the form $\trace = [\event_1] \tracecat \trace' \tracecat [\event_2]$ with $\trace' \in \ftdom$, 
	and $\eventdep(\event_1, \event_2, \trace, c, D)$,
then for every $\event \in \trace'$, if it either doesn't have the \emph{Event May-Dependency} relations on $\event_1$,
or $\pi_1(\event) \notin \kw{FV}(\pi_1(\event_2)) $,
then 
$\pi_1(\event_1) \in \kw{FV}(\pi_1(\event_2))$, and $ {\pi_2(\event_1)} = \llabel(\trace)$
%
If $\event_2 \in \eventset^{\test}$, 
then for every $\event \in \trace'$, if it either doesn't have the \emph{Event May-Dependency} relations on $\event_1$,
or $\pi_1(\event) \notin \kw{FV}(\pi_1(\event_2)) $,
then 
$\pi_1(\event_1) \in \kw{FV}(\pi_1(\event_2))$, and $ {\pi_2(\event_1)} = \llabel(\trace)$
%
\[
\begin{array}{l}
	\forall D \in \dbdom , c \in \cdom, \trace \in \ftdom \sthat  \forall \event_1,\in \eventset^{\asn}, \event_2 \in \eventset^{\test} \sthat 
	 \exists \trace' \in \ftdom \sthat
	 \\ \quad 
	 \trace = [\event_1] \tracecat \trace' \tracecat [\event_2]
	 \land
	\eventdep(\event_1, \event_2, \trace, c, D) 
	\\ \quad 
	\land 
	\Big( \nexists \event \in \trace' \sthat  
	\pi_1(\event) \in \kw{FV}(\pi_1(\event_2)) \land
	\eventdep(\event_1, \event, \trace[\event_1:\event], c, D)
	\Big) 
	\\ \quad 
	\implies 
	\pi_1(\event_1) \in \kw{FV}(\pi_1(\event_2)) \land {\pi_2(\event_1)} = \llabel(\trace)
\end{array}
\]
\end{lem}
%
\begin{proof}
Taking arbitrary $D \in \dbdom , c \in \cdom$, and an assignment events $\event_1 \in \eventset^{\asn}$ and another event 
$\event_2\in \eventset$.
\\
Without loss of generalization, 
taking arbitrary trace has the form $\trace = [\event_1, \ldots, \event_2]$ 
% arbitrary $\trace_2 \in \ftdom$,
%  then we know 
 and $\trace_2 \in \ftdom$ such that
%   $ $ and
 \[
	\trace = [\event_1] \tracecat \trace_2 \tracecat [\event_2]
	\land
	\eventdep(\event_1, \event_2, \trace, c, D)
	\]
By the definition of $\eventdep(\event_1, \event_2, \trace, c, D)$, 
taking $ \event_1' \in \eventset^{\asn},
\trace_2' \in \ftdom, c_1, c_2 \in \cdom$ and $\event_2' \in \eventset^{\test}$
 as the events, traces and commands satisfying the definition,
 we have following two executions:
% \[
%   \exists \event_1', \event_2' \in \eventset^{\asn},
%   \trace_2' \in \ftdom, c_1, c_2 \in \cdom \sthat 
% \]
%
\[
\begin{array}{l}
\config{c, \trace_0} \rightarrow^{*}
\config{c_1, \trace_1 \tracecat [\event_1]} \rightarrow^{*} \config{c_2, \trace_1 \tracecat [\event_1] \tracecat \trace_2 \tracecat [\event_2]} 
\\ \quad
% \land
\config{c_1, \trace_1 \tracecat [\event_1']} \rightarrow^{*} \config{c_2, \trace_1 \tracecat [\event_1'] \tracecat \trace_2' \tracecat [\event_2']} 
\end{array}
\]
%
Taking arbitrary $\event_z \in \trace_2$, we know 
$\neg \eventdep(\event_1, \event, \trace[\event_1:\event_z], c, D)
\lor  \pi_1(\event_z) \notin \kw{FV}(\expr_2 / \qexpr_2)$.
\\
Then by the same proof in Lemma~\ref{lem:inv_indepevents_asn}, and applying the Inversion Lemma~\ref{lem:inv_expr} and \ref{lem:inv_expr_gnl} of the boolean expression case,
we have:
\[ 
	\pi_1(\event_1) \in \kw{FV}(\pi_1(\event_2)) \land {\pi_2(\event_1)} = \llabel(\trace)
	\]
	This case is proved.
\end{proof}

\begin{lem}[Control Dependency Inversion]
	\label{lem:ctldep_inv}
	For every $c \in \cdom$, $D \in \dbdom, \trace \in \ftdom$ and two assignment events $\event_1, \event_2 \in \eventset^{\asn} $, if $\eventdep(\event_1, \event_2, c, \trace, D)$ and
	satisfying the Equation~\ref{eq:eventdep}(c) in Definition~\ref{def:event_dep},
	% they are in the second case of the \emph{Event May-Dependency} relation from Definition~\ref{def:event_dep},
	% $\eventdep(\event, \event, c, \trace, D)$ as Eq.~\ref{eq:ctlflowsto_inv},
	then for all  $z \in \kw{FV}(\pi_1(\event_b))$ there exists a label $i \in \mathbb{N}$ such that 
	$\flowsto(z^i, \pi_1(\event)^{\pi_2(\event)}, c)$
	\begin{equation}
		\label{eq:ctlflowsto_inv},		
		\begin{array}{l}
			\forall D \in \dbdom , c \in \cdom, \trace \in \ftdom,
			\event_1, \event_2 \in \eventset^{\asn} \sthat  
			\\ 
			\exists \vtrace_0,
			\vtrace_1, \vtrace', \vtrace_3, \vtrace_3' \in \ftdom,\event_1' \in \eventset^{\asn}, {c}_1, {c}_2  \in \cdom, 
			\event_b \in \eventset^{\test},
			\trace_{ih} \in \ftdom \sthat  
		\trace = [\event_1] \tracecat \trace_{ih} \tracecat [\event_2]
		\\ \quad \land	  
			  \config{{c}, \vtrace_0} \rightarrow^{*} 
				\config{{c}_1, \vtrace_1 \tracecat [\event_1]}  \rightarrow^{*} 
				\config{c_2,  \vtrace_1 \tracecat [\event_1] \tracecat \trace \tracecat [\event_b] \tracecat  \trace_3} 
			  \\ \qquad \land
			  \config{{c}_1, \vtrace_1 \tracecat [\event_1']}  \rightarrow^{*} 
			  \config{c_2,  \vtrace_1 \tracecat [\event_1'] \tracecat \trace' \tracecat [(\neg \event_b)] \tracecat \trace_3'} 
			  \\ \qquad \land
			\tlabel({\trace_3}) \cap \tlabel({\trace_3'}) = \emptyset
			   \land \vcounter(\trace') ~  \pi_2(\event_b) = \vcounter(\trace) ~  \pi_2(\event_b)
				\land \event_2 \eventin \trace_3
			  \land \event_2 \not\eventin \trace_3'
		\\ \quad \implies	
		\forall z \in \kw{FV}(\pi_1(\event_b)) \sthat  
		\exists l \in \mathbb{N} \sthat  
		\flowsto(z^l, \pi_1(\event_2)^{\pi_2(\event_2)},c)
	\end{array}
\end{equation}
	\end{lem}
	Proof Summary:
	\\
	Proving by using the Inversion Lemmas~\ref{lem:inv_expr}, \ref{lem:inv_expr_gnl}, 
	\ref{lem:inv_event}, and \ref{lem:inv_live}, and the \emph{Event May-Dependency} definition of the second case.
%
	\begin{proof}
		Take arbitrary $D \in \dbdom , c \in \cdom, \trace \in \ftdom,
		\event_1, \event_2 \in \eventset^{\asn} $,
%
let $\vtrace_0,
\vtrace_1, \vtrace', \vtrace_3, \vtrace_3' \in \ftdom, 
\event_2' \in \eventset, \event_1' \in \eventset^{\asn}, \event_b, {c}_1, {c}_2$ be the traces, 
events and commands satisfying the executions,
by Inversion Lemma~\ref{lem:inv_event} on 
$\event_2$, and $\event_b$,
we have the following instance of the first execution in Eq.~\ref{eq:ctlflowsto_inv},
 %
%
% Let $\event_{ih} = (b, l_b, n_b, v_b)$, by Eq.~\ref{eq:ctldep_inv1} and {Inversion Lemma~\ref{lem:inv_test}}, we have:
\begin{equation}
% \label{eq:ctldep_inv1}
  \begin{array}{l}   
\config{{c}, \vtrace_0} \rightarrow^{*} 
% \config{[\assign{{x}_1}{\expr_1 / \query(\qexpr_1)}]^{l_1} ; {c}_1, \vtrace_1}  
% \rightarrow^{\rname{assn/query}}
%  \config{c_1, \vtrace_1 \tracecat [\event_1]} 
%  \\ \qquad 
%  \rightarrow^{*} 
  \config{\eif ([b]^{l_b}, c_t, c_f) / \ewhile [b]^{l_b} \edo (c_w);{c}_3', 
  \vtrace_1 \tracecat [\event_1] \tracecat \trace} 
  \\
  \qquad 
   \rightarrow^{\rname{if-b / while-b}} 
  \config{(c_t;c_3' / c_f;c_3') /(c_w; \ewhile [b]^{l_b} \edo (c_w);c_3'/\clabel{\eskip}^{l'}; c_3'), 
  \trace_1 \tracecat [\event_1] \tracecat \trace \tracecat [\event_b]} 
  \\
  \qquad  \rightarrow^{*} 
  \config{[\assign{{x}_2}{\expr_2 / \query(\qexpr_2)}]^{l_2}; c_{3b}', 
  \trace_1 \tracecat [\event_1] \tracecat \trace \tracecat [\event_b] \tracecat  \trace_{3a}}
  \\ \qquad \rightarrow^{\rname{assn/query}}
  \config{ c_{3b}', 
  \trace_1 \tracecat [\event_1] \tracecat \trace \tracecat [\event_b] \tracecat  \trace_{3a} \tracecat [\event_2]}
  \\ \qquad
  \rightarrow^{*} 
  \config{c_3, 
  \trace_1 \tracecat [\event_1] \tracecat \trace \tracecat [\event_b] \tracecat  \trace_{3a} \tracecat [\event_2] \tracecat \trace_{3b}}
  % 
\end{array}
\end{equation}
%  %
, where $\trace_3 = \trace_{3a} \tracecat [\event_2] \tracecat \trace_{3b}$,
% $x_1 = \pi_1(\event_1)$, $l_1 = {\pi_2(\event_1)}$, 
$x_2 = \pi_1(\event_2)$, $l_2 = \pi_2(\event_2)$, 
and $\eif ([b]^{l_b}, c_t, c_f) / \ewhile [b]^{l_b} \edo (c_w)$ 
is the conditional command of the assignment commands associated to the 
$\event_b$ from Inversion Lemma~\ref{lem:inv_event} of testing event.
\\
% \[\begin{array}{l}   
% 	  \config{\eif ([b]^{l_b}, c_t, c_f) / \ewhile [b]^{l_b} \edo (c_w);{c}_3', 
% 	  \vtrace_1 \tracecat [\event_1] \tracecat \trace} 
% 	  \\
% 	  \qquad 
% 	   \rightarrow^{\rname{if-b / while-b}} 
% 	  \config{(c_t;c_3' / c_f;c_3') /(c_3' / c_w; \ewhile [b]^{l_b} \edo (c_w);{c}_3'), 
% 	  \trace_1 \tracecat [\event_1] \tracecat \trace \tracecat [\event_b]} 
% 	  % 
% 	\end{array}
% 	\]
% $\config{\eif ([b]^{l_b}, c_t, c_f) / \ewhile [b]^{l_b} \edo (c_w);{c}_3', 
% \vtrace_1 \tracecat [\event_1] \tracecat \trace} 
%  \rightarrow^{\rname{if-b / while-b}} 
% \config{(c_t;c_3' / c_f;c_3') /(c_3' / c_w; \ewhile [b]^{l_b} \edo (c_w);{c}_3'), 
% \trace_1 \tracecat [\event_1] \tracecat \trace \tracecat [\event_b]} 
% $
The notation $(c_t;c_3' / c_f;c_3') /(c_w; \ewhile [b]^{l_b} \edo (c_w);c_3' / \clabel{\eskip}^{l'}; c_3')$ represents:
\\
In case of $\eif ([b]^{l_b}, c_t, c_f)$, if $\pi_3(\event_b) = \etrue$, we have the evaluation:
$$
\config{\eif ([b]^{l_b}, c_t, c_f) ;{c}_3', 
\vtrace_1 \tracecat [\event_1] \tracecat \trace} 
 \rightarrow^{\rname{if-b}} 
\config{c_t;c_3' 
\trace_1 \tracecat [\event_1] \tracecat \trace \tracecat [\event_b]} 
$$
%
The same for case of $\eif ([b]^{l_b}, c_t, c_f)$ with $\pi_3(\event_b) = \efalse$,
and case of $\ewhile [b]^{l_b} \edo (c_w)$ with $\pi_3(\event_b) = \etrue$ and $\pi_3(\event_b) = \efalse$.
%
\\
By the command label consistency,
we also have the instance of second execution as follows:
\begin{equation}
\label{eq:ctldep_inv2}
\begin{array}{l}   
  \config{{c}, \vtrace_0} \rightarrow^{*} 
%   \config{[\assign{{x}_1}{\expr_1 / \query(\qexpr_1)}]^{l_1} ; {c}_1, \vtrace_1}  
%   \rightarrow^{\rname{\rname{assn/query}}}
%    \config{c_1, \vtrace_1 \tracecat [\event_1]} 
%    \\
%     \qquad \rightarrow^{*} 
    \config{\eif ([b]^{l_b}, c_t, c_f) / \ewhile [b]^{l_b} \edo (c_w);{c}_3', 
    \vtrace_1 \tracecat [\event_1] \tracecat \trace'} 
    \\
    \qquad 
     \rightarrow^{\rname{if-b / while-b}} 
    \config{(c_f;c_3' / c_t;c_3') /(\clabel{\eskip}^{l'}; c_3' / c_w; \ewhile [b]^{l_b} \edo (c_w);{c}_3' ), 
	\\ \qquad \qquad \qquad \qquad
	\trace_1 \tracecat [\event_1]  \tracecat \trace' \tracecat [\neg \event_b]} 
    \\
    \qquad   \rightarrow^{*} 
    \config{c_3, 
    \trace_1 \tracecat [\event_1]  \tracecat \trace' \tracecat [\neg \event_b] \tracecat  \trace_3'}
    % 
  \end{array}
\end{equation}
%
By the label consistency, and $\tlabel({\trace_3}) \cap \tlabel({\trace_3'}) = \emptyset$, 
% i.e., 
we know $\trace_3$ and $\trace_3'$ doesn't contain any event of evaluating the commands in $c_3'$.
Otherwise, $\tlabel({\trace_3}) \cap \tlabel({\trace_3'}) \neq \emptyset$, which is a contradiction.
\\
Since $\trace_3= \trace_{3a} \tracecat [\event_2] \tracecat \trace_{3b} $, we know $\event_2$ doesn't comes from evaluating
of $c_3'$, i.e.,:
\\
In the case of $\eif ([b]^{l_b}, c_t, c_f)$, $\event_2$ comes from the evaluation of $c_t$ or $c_f$,
Then we know
\\
 $[\assign{{x}_2}{\expr_2 / \query(\qexpr_2)}]^{l_2} \in_c c_t 
 \lor [\assign{{x}_2}{\expr_2 / \query(\qexpr_2)}]^{l_2} \in_c c_f$;
\\
and in the case of $\ewhile [b]^{l_b} \edo (c_w)$, $\event_2$ comes from the evaluation of $c_w$,
i.e., $[\assign{{x}_2}{\expr_2 / \query(\qexpr_2)}]^{l_2} \in_c c_w$.
\\
In both of the two cases, we know $\forall z \in \kw{FV}(\pi_1(\event_b)) $ there is a label $l \in \mathbb{N}$ for this variable,
and by the $\flowsto$ definition, 
$\flowsto(z^l, \pi_1(\event_2)^{\pi_2(\event_2)},c)$.
\\
This lemma is proved.
	\end{proof}
%
\input{appendices/thm-adaptivity/lem_basicinversion}
%
\begin{lem}[While Loop Inversion]
	\label{lem:inv_while}
	For every $\trace, \trace' \in \ftdom, c, c_1, c_2 \in \cdom$ 
	if $ \config{c, \trace} \rightarrow^* \config{c_1; c_2, \trace'}$ and 
	$c_1 \in_c c_2$, 
	then there must exist a $\ewhile$ command in $c_2$ and $c_1$ must shows up in the body of that $\ewhile$ command,
	 i.e., $\exists l \in \mathbb{N}, b \in \mathcal{B}, c_w \in \cdom \sthat  
	(\ewhile [b]^l \edo (c_w)) \in_c c_2 \land c_1 \in_c c_w$.
	%
	\[
	\begin{array}{l}
	\forall \trace, \trace' \in \ftdom, c, c_1, c_2 \in \cdom \sthat 
		\\ \quad
		\config{c, \trace} \rightarrow^* \config{c_1; c_2, \trace'}
		\implies
		\\ \qquad
		(c_1 \in_c c_2
		\land
		\exists l \in \mathbb{N}, b \in \mathcal{B}, c_w \in \cdom \sthat  
		(\ewhile [b]^l \edo (c_w)) \in_c c_2 \land c_1 \in_c c_w)
	\end{array}
	\]
	\end{lem}	
	Proof Summary: trivially by induction on $c$ and enumerate all operational semantic rules.
\begin{proof}
	Take arbitrary $\trace \in \ftdom$, by induction on $c$, we have following cases:
		\caseL{$c = [\assign{x}{\expr}]^l$}
		By the evaluation rule $\rname{assn}$, we have
		$
		{
		\config{[\assign{{x}}{\aexpr}]^{l},  \trace } 
		\xrightarrow{} 
		\config{\clabel{\eskip}^{l'}, \trace \tracecat [({x}, l, v) ]}
		}$.
		\\
		Since there doesn't exist $c_1, c_2 \in \cdom$ satisfying $\eskip = c_1; c_2$, this theorem is vacuously true.
		\caseL{$c = [\assign{x}{\query(\qexpr)}]^l$}
		By the evaluation rule $\rname{query}$, we have
		$
		{
		\config{[\assign{{x}}{\query(\qexpr)}]^{l},  \trace } 
		\xrightarrow{} 
		\config{\clabel{\eskip}^{l'}, \trace \tracecat [({x}, l, \qval, v) ]}
		}$.
		\\
		Since there doesn't exist $c_1, c_2 \in \cdom$ satisfying $\eskip = c_1; c_2$, this theorem is vacuously true.
		\caseL{$c = \eif([b]^{l}, c_1, c_2)$}
		By the evaluation rule $\rname{query}$ and $\rname{if-f}$, and the label consistency, we know:
		\\
		for all possible $c_{t1}$ and $c_{t2}$ 
		such that $c_t$ has the form $c_t = c_{t1};c_{t2}$;
		\\
		all possible $c_{f1}$ and $c_{f2}$ 
		such that $c_f$ has the form $c_f = c_{f1};c_{f2}$,
		\\
		$c_{t1} \notin c_{t1}$ and $c_{f1} \notin c_{f2}$.
		\\
		Then this theorem is vacuously true.
		\caseL{$c = c_{s1};c_{s2}$}
		By label consistency, we know for every $c_1' \in_c c_{s1}$, $c_1' \notin c_{s2}$.
		\\
		Then by the induction hypothesis on $c_{s1}$ and $c_{s2}$ separately, we have this case proved.
		\caseL{$\ewhile [b]^{l} \edo c$}
		By rule $\rname{while-t}$, we have:
		\[
			\config{{\ewhile [b]^{l} \edo (c_w), \trace}}
			\xrightarrow{} 
			\config{{
			c_w; \ewhile [b]^{l} \edo (c_w),
			\trace \tracecat [\event]}}
		\]
		%
		If $c_w$ is a sequence command,
		let $c_1 = c_{w1}$ be the any possible command in this sequence, for all possible $c_{w1}$ and $c_{w2}$ 
		such that $c_w$ has the form $c_w = c_{w1};c_{w2}$.
		\\
		Then we have $c_2 = c_{w2};\ewhile [b]^{l} \edo (c_w)$ and $c_1 \in_c c_2$.
		\\
		And we also have the existence of $l = l_b, b$ and $c_w$, and $\ewhile [b]^{l} \edo (c_w) \in_c c_2$ and  $c_1 \in c_w$.
		\\
		If $c_w$ isn't a sequence command, let $c_1 = c_w$, then we have $c_2 = \ewhile [b]^{l} \edo (c_w)$ 
		and $c_1 \in_c c_2$.
		\\
		And we also have the existence of $l = l_b, b$ and $c_w$, and $\ewhile [b]^{l} \edo (c_w) \in_c c_2$ and  $c_1 \in c_w$.
		\\
		This case is proved.
		\\
		By the evaluation rule $\rname{while-f}$, we have
		$$
		{
			\config{{\ewhile [b]^{l} \edo (c_w), \trace}}
			\xrightarrow{} 
			\config{{
			\clabel{\eskip}^l ,
			\trace \tracecat [(b, l, \efalse)]}}.
		}$$
		\\
		Since there doesn't exist $c_1, c_2 \in \cdom$ satisfying $\eskip = c_1; c_2$, this theorem is vacuously true.
	\end{proof}
%
%
\begin{lem}[Only $\eskip$ Command doesn't Produce Event].
	\label{lem:inv_skip}
	For all trace $\trace\in \ftdom$, and $c, c' \in \cdom$,  
	$\config{c, \trace} \rightarrow \config{c', \trace}$ if and only if $c = \clabel{\eskip}^{l'};c'$. 
	\[
		\forall \trace\in \ftdom, c, c' \in \cdom \sthat 
		\config{c, \trace} \rightarrow \config{c', \trace}
		\Leftrightarrow 
		c = \clabel{\eskip}^{l'};c'
	% \footnote{$(\clabel{\eskip}^{l'};){}^*$ denotes a sequence command only composed of $\clabel{\eskip}^{l'}$ commands.}
	\]
	\end{lem}
\begin{proof}
	Proved trivially by induction on $c$ and enumerate all operational semantic rules.
\end{proof}
%
% \todo{Event Dependency Transitivity}
\begin{lem}
	\label{lem:valdep_trans}
(Event Dependency Transitivity)
For every $D \in \dbdom , c \in \cdom, \trace \in \ftdom$, and $\event_1, \event_2, \event_3 \in \eventset^{\asn}, \trace_{12}, \trace_{23} \in \ftdom$,
if $\eventdep(\event_1, \event_2, \trace_{12}, c, D)$
and $\eventdep(\event_2, \event_3, \trace_{23}, c, D) $,
then there is an event dependency relation,
$\eventdep(\event_1, \event_3, \trace_{12}\tracecat\trace_{23}, c, D)$.
  % An event $\event_2 \in \eventset^{\asn}$ is in the \emph{may-dependency} relation with another
  % event $\event_1 \in \eventset^{\asn}$ in a program ${c}$ with a hidden database $D$, denoted as $\eventdep(\event_1, \event_2, c, D)$,
  % if and only if
  \[
	  \begin{array}{l}
  \forall D \in \dbdom , c \in \cdom, \event_1, \event_2, \event_3 \in \eventset^{\asn}, \trace_{12}, \trace_{23} \in \ftdom \sthat  
  \\ \quad
  \eventdep(\event_1, \event_2, \trace_{12}, c, D) 
  \land
  \eventdep(\event_2, \event_3, \trace_{23}, c, D) 
  \implies
  \eventdep(\event_1, \event_3, \trace_{12}\tracecat\trace_{23}, c, D)
	  \end{array}
  \]
\end{lem}
%
%

%
\begin{lem}[\emph{Variable May-Dependency} Transitivity]
	\label{lem:vardep_trans}
For every $c \in \cdom, x^i, y^j, z^l \in \lvar(c)$, 
if $\vardep(x^i, y^j, c)$ and 
$\vardep(y^j, z^l, c)$, then $\vardep(x^i, z^l, c)$.
	\[
	\forall c \in \cdom, x^i, y^j, z^l \in \lvar(c) \sthat  
	\vardep(x^i, y^j, c) 
	\land
	\vardep(y^j, z^l, c) 
	\implies
	\vardep(x^i, z^l, c)
	\]
  \end{lem}
  %
%
